%\documentclass{...}

\input{./def/yf-formatting}

\usepackage{amsfonts, amsmath, amssymb, amsthm}
\usepackage{comment}
\usepackage{graphicx}
\usepackage{ifthen}
\usepackage{latexsym}
%\usepackage{times}
\usepackage[normalem]{ulem}

\input{./def/yf-def}

\def\dom{\prec}
\def\T{\mathcal{T}}

\newboolean{solver}\setboolean{solver}{true}
%\newboolean{solver}\setboolean{solver}{false}
\ifthenelse{\boolean{solver}}{\includecomment{sol}}{\excludecomment{sol}}

\begin{document}

\section*{CE7456: Quiz 2}

\noindent Name:  \hspace{50mm}
Student ID:

% \begin{center}
%     \uline{Classroom Discussion}
% \end{center}

\extraspacing


\extraspacing {\bf Problem 1 (25 marks).} Let $P$ be a set of $n$ points in $\real^2$. Given a point $q \in \real^2$, a {\em nearest neighbor} (NN) query returns a point in $P$ that has the smallest Euclidean distance to $q$. The goal of the NN search problem is to preprocess $P$ into a data structure that can answer NN queries efficiently. It is known that we can build a structure in $O(n \log n)$ time to answer NN queries in $O(\log n)$ time.

\vgap

Prove: The NN search problem is $(\log_2 n, \log_2 n)$ spectrum indexable.

\pagebreak

\extraspacing {\bf Problem 2 (25 marks).} Consider the set $S = \set{10, 30, 35, 40, 70, 75, 85, 90}$. Answer the following questions:

\myitems{
    \item Show the ``learned index'' discussed in class on $S$.

    \item Describe how to find the predecessor of $q = 50$ in $S$ using the index.
}

\pagebreak

\extraspacing {\bf Problem 3 (25 marks).} A 6-cycle is an (undirected) cycle with 6 vertices. Answer the following question:

\myitems{
    \item What is the fractional edge cover number $\rho$ of 6-cycle?

    \item Given an undirected graph $G = (V, E)$ with $m$ edges, describe how to find all the 6-cycles in $G$ in $O(m^\rho)$ time.

    \vgap

    Note: ``apply the algorithm taught in the class'' is not an acceptable answer. You need to describe the detailed steps.
}

\pagebreak

 \extraspacing {\bf Problem 4 (25 marks).} Consider a join $\mathcal{Q}$ with four relations $R_1(ABC), R_2(ABD), R_3(ACD)$, and $R_4(BCD)$. Here, the letters in parentheses illustrate relation schemas (e.g., the schema of $R_1$ is $\set{A, B, C}$).

 \vgap

 Prove: If each relation has $n$ tuples, then the AGM bound of $\mathcal{Q}$ is $O(n^{4/3})$.
\end{document}
